\subsection{Cardinality of Poincaré groups}

PROPOSITION 9. --- Let X be a loop-disentanglable topological space and let W be a base of the topology of X. For every point x of X, the cardinality of the group \(\pi_{1}(X,x)\) is bounded above by sup(Card(W),Card(N)). In particular, the Poincaré group of a metric space of countable type that is loop-disentanglable is countable.

Indeed, the space \(\lambda_{x}(\mathrm{X})\) equipped with the endpoint map is a connected covering of X whose fiber at x is \(\pi_{1}(\mathrm{X}, x)\) . The assertion then follows from I, p. 40, th. 3.

Lemma 4. --- Let X be a Baire space, let G be a topological group acting continuously on X, and let B be an approachable subset of X that is not meager. The set of points \(g \in G\) such that \(B \cap gB \neq \emptyset\) is a neighborhood of the identity element of G.

Since B is approachable, \(\complement B\) is also approachable, because the collection of approachable subsets of X is a \(\sigma\)-algebra. There therefore exists an open subset U of X such that \(U \cap \complement B\) and \(B \cap \complement U\) are meager in X.

Let V be a neighborhood of the identity element of G and W a nonempty open subset of X contained in U such that \(V \cdot W \subset U\) . For all \(g \in V\) , \(U \cap gU\) is not empty.

Let \(g \in G\) such that \(B \cap gB\) is empty. The relations

\[
\begin{array}{l} \mathrm{U} \cap g \mathrm{U} = (\mathrm{U} \cap g \mathrm{U}) \cap \mathbb {C} (\mathrm{B} \cap g \mathrm{B}) \\ \qquad = (\mathrm{U} \cap g \mathrm{U}) \cap (\mathbb {C} \mathrm{B} \cup g \mathbb {C} \mathrm{B}) \\ \qquad = (\mathrm{U} \cap g \mathrm{U} \cap \mathbb {C} \mathrm{B}) \cup (\mathrm{U} \cap g \mathrm{U} \cap g \mathbb {C} \mathrm{B}) \\ \qquad \subset (\mathrm{U} \cap \mathbb {C} \mathrm{B}) \cup g (\mathrm{U} \cap \mathbb {C} \mathrm{B}) \end{array}
\]

imply that U ∩ gU is meager in X. Since it is an open subset and X is a Baire space, it is empty. We therefore have \(g \notin V\) , whence the lemma.

THEOREM 2 (Shelah \(^{(2)}\) ). --- Let X be a connected and locally path-connected Polish space and let x be a point of X. If X is not loop-disentanglable, the group \(\pi_{1}(X,x)\) has the cardinality of the continuum.

Let d be a metric defining the topology of X for which it is complete. Suppose that X is not loop-disentanglable. Then there exists a point \(a \in X\) and, for every integer \(n \geqslant 0\) , a loop \(c_{n}\) in X whose image has diameter \(\leqslant 2^{-n}\) and whose class in \(\pi_{1}(X, a)\) is not trivial.

Let K denote the set \(\{0,1\}^{\mathbf{N}}\) and endow it with the product topology, the space \(\{0,1\}\) equipped with the discrete topology. For every element \(\varepsilon = (\varepsilon_n)\) of K, let \(c_{\varepsilon}\) be the map from I to X defined by \(c_{\varepsilon}(0) = a\) and, for \(2^{-n - 1} \leqslant t \leqslant 2^{-n}\) , \(c_{\varepsilon}(t) = c_n(2^{n + 1}t - 1)\) if \(\varepsilon_n = 1\) and \(c_{\varepsilon}(t) = a\) otherwise. We have \(c_{\varepsilon}(0) = c_{\varepsilon}(1) = a\) . The map \(c_{\varepsilon}\) is continuous at every point \(t \in ]0,1]\) . It is also so at 0 because \(d(a,c_{\varepsilon}(t)) \leqslant 2^{-n}\) if \(t \in [0,2^{-n}] \) . The map \(c_{\varepsilon}\) is therefore a loop in X at \(a\) .

If \(\varepsilon\) and \(\varepsilon'\) are elements of K such that \((\varepsilon_0, \ldots, \varepsilon_n) = (\varepsilon_0', \ldots, \varepsilon_n')\) , then \(c_{\varepsilon}(t) = c_{\varepsilon'}(t)\) for all \(t \in [2^{-n-1}, 1]\) and \(d(c_{\varepsilon}(t), c_{\varepsilon'}(t)) \leqslant d(a, c_{\varepsilon}(t)) + d(a, c_{\varepsilon'}(t)) \leqslant 2^{-n}\) if \(t \in [0, 2^{-n-1}]\) . It follows that the map \(\varepsilon \mapsto c_{\varepsilon}\) from K into the space \(\Omega_a(X)\) is continuous, when the space \(\Omega_a(X)\) is equipped with the topology of compact convergence.

Denote by \(\Gamma \subset \mathrm{K} \times \mathrm{K}\) the set of pairs \((\varepsilon, \varepsilon')\) such that \(c_{\varepsilon}\) is strictly homotopic to \(c_{\varepsilon'}\) . This is the graph of an equivalence relation R on K.

Lemma 5. --- The set \(\Gamma\) is a meager subset of K × K.

Let Z denote the space \(K \times K \times \mathcal{C}_{c}(\mathbf{I} \times \mathbf{I}; X)\) . The topology of the space \(\mathcal{C}_{c}(\mathbf{I} \times \mathbf{I}; X)\) is defined by the metric \(\delta\) given by \(\delta(h, h') = \sup_{u \in \mathbf{I} \times \mathbf{I}} d(h(u), h'(u))\) and it is complete for this metric (TG, X, p. 20, corollary and TG, X, p. 9, cor. 1); by TG, X, p. 24, th. 1, this topology is of countable type. The space \(\mathcal{C}_{c}(\mathbf{I} \times \mathbf{I}; X)\) is therefore a Polish space. The same holds for the space Z, since K is a Polish space (TG, IX, p. 57, prop. 1).

Let H be the subset of Z consisting of the triples \((\varepsilon,\varepsilon^{\prime},h)\) such that h is a strict homotopy connecting \(c_{\varepsilon}\) to \(c_{\varepsilon^{\prime}}\) . The maps from Z into \(X^{2}\) given by \(a_{t}\colon(\varepsilon,\varepsilon^{\prime},h)\mapsto(c_{\varepsilon}(t),h(t,0))\) , \(b_{t}\colon(\varepsilon,\varepsilon^{\prime},h)\mapsto(c_{\varepsilon^{\prime}}(t),h(t,1))\) and \(c_{t}\colon(\varepsilon,\varepsilon^{\prime},h)\mapsto(h(0,t),h(1,t))\) are continuous, for every \(t\in I\) , because the map \(\varepsilon\mapsto c_{\varepsilon}\) from K to \(\Omega_{a}(X)\) is continuous, as are the mappings \(h\mapsto h(s,t)\) of \(\mathcal{C}_{c}(\mathbf{I}\times\mathbf{I};\mathbf{X})\) into X. By definition, H is the intersection of the sets \(a_{t}^{-1}(\Delta_{\mathrm{X}})\) , \(b_{t}^{-1}(\Delta_{\mathrm{X}})\) and \(c_{t}^{-1}((a,a))\) , for \(t\in I\) , where \(\Delta_{X}\) denotes the diagonal of X. This proves that H is a closed subset of Z.

Consequently, H is a Polish space. Let \(p \colon \mathrm{Z} \to \mathrm{K} \times \mathrm{K}\) be the canonical projection; we have \(\Gamma = p(\mathrm{H})\) by definition. Since \(\mathrm{K} \times \mathrm{K}\) is separated, \(\Gamma\) is a Souslin subset of \(\mathrm{K} \times \mathrm{K}\) (TG, IX, p. 59, def. 2). According to IV, p. 362, prop. 7, this is an approachable subset of \(\mathrm{K} \times \mathrm{K}\) .

Suppose that \(\Gamma\) is not meager. The space K × K is a compact topological space, hence a Baire space (TG, IX, p. 55, th. 1). Endow the group \(G_{0}\) of permutations of the set \(\{0,1\}\) with the discrete topology, and let the product topological group \(G = G_{0}^{\mathbf{N}}\) act diagonally on \(K = \{0,1\}^{\mathbf{N}}\) . The group G then acts continuously on K × K via the map (g, (x,y)) ↦ (x, g · y). By Lemma 4, the set V of elements \(g \in G\) such that \(\Gamma \cap g \cdot \Gamma \neq \varnothing\) is a neighborhood of the identity element of G.

Let \(g \in V\) ; let \((\varepsilon, \varepsilon') \in \Gamma \cap g \cdot \Gamma\) . Since \((\varepsilon, \varepsilon') \in \Gamma\) , we have \(\varepsilon\,R\,\varepsilon'\) ; since \((\varepsilon, \varepsilon') \in g \cdot \Gamma\) , \(g^{-1} \cdot (\varepsilon, \varepsilon') = (\varepsilon, g^{-1}\cdot\varepsilon') \in \Gamma\) , so that \(\varepsilon\,R\,g^{-1}\cdot\varepsilon'\) . We have thus shown that, for every \(g \in V\) , there exists \(\varepsilon \in K\) such that \(\varepsilon\) and \(g\varepsilon\) are equivalent for R.

For \(m \in \mathbf{N}\) , denote by \(\tau_m\) the element of G all of whose terms are equal to \(e\) except the one with index \(m\) which is equal to the nontrivial element \(\tau\) of \(G_0\) . There exists an integer \(m\) such that \(\tau_m\) belongs to V; then choose \(\varepsilon \in K\) such that \(\varepsilon\) and \(\varepsilon' = \tau_m \cdot \varepsilon\) are equivalent for R. This implies that the loops \(c_{\varepsilon}\) and \(c_{\varepsilon'}\) are strictly homotopic. By construction, these loops coincide on the intervals \([0, 2^{-m-1}]\) and \([2^{-m}, 1]\) ; on the interval \([2^{-m-1}, 2^{-m}]\) , one is the constant map with image a and the other is the map \(t \mapsto c_{m}(2^{m+1}t - 1)\) . It follows that \(c_{m}\) is strictly homotopic to the constant loop with image \(\{a\}\) , whence a contradiction. Lemma 5 is thus proved.

We now finish the proof of Theorem 2. By Prop. 8, there exists a continuous injective map \(\gamma\) of \(\{0,1\}^{\mathbf{N}}\) into K whose image meets every equivalence class under R in at most one point. If \(k\) and \(k'\) are distinct elements of \(\{0,1\}^{\mathbf{N}}\) , the loops \(c_{\gamma(k)}\) and \(c_{\gamma(k')}\) are not strictly homotopic in X and the map \(\{0,1\}^{\mathbf{N}} \to \pi_1(\mathrm{X},a)\) given by \(k \mapsto [c_{\gamma(k)}]\) is injective. In particular, \(\operatorname{Card}(\pi_1(\mathrm{X},a)) \geqslant \operatorname{Card}(\{0,1\}^{\mathbf{N}}) = \operatorname{Card}(\mathfrak{P}(\mathbf{N}))\) . Since X is a metrizable topological space of countable type, so is \(\Omega_a(\mathrm{X})\) (TG, X, p. 24, th. 1). Consequently, \(\Omega_a(\mathrm{X})\) is homeomorphic to a subspace of \([0,1]^{\mathbf{N}}\) (TG, IX, p. 18, prop. 12) and

\[
\begin{array}{r l} & {\mathrm{Card} (\Omega_ {a} (\mathbf {X})) \leqslant \mathrm{Card} ([ 0, 1 ] ^ {\mathbf {N}}) = \mathrm{Card} (\mathfrak {P} (\mathbf {N}) ^ {\mathbf {N}})} \\ & {\qquad = \mathrm{Card} (\mathfrak {P} (\mathbf {N} \times \mathbf {N})) = \mathrm{Card} (\mathfrak {P} (\mathbf {N})).} \end{array}
\]

A fortiori, \(\operatorname{Card}(\pi_{1}(\mathrm{X},a))\leqslant\operatorname{Card}(\mathfrak{P}(\mathbf{N}))\) . It then follows from cor. 2 of E, III, p. 25, that \(\pi_{1}(\mathrm{X},a)\) has the cardinality of the continuum, as was to be proved.

Example. --- Let P be the topological space which is the union of the circles with center \((2/n,0)\) passing through the origin of the plane \(R^{2}\) , for \(n \geqslant 1\) (III, p. 336, exerc. 6). The Poincaré group of P has the cardinality of the continuum (III, p. 338, exerc. 9).

